\section{Contract DAGs and reverse demand propagation}
\label{sec:dag}

\subsection{Why local composition is not enough}

A large formalization is rarely a single loop.  It is a directed acyclic graph
of constructions: coefficient approximation feeds root isolation and an error
estimate; both feed a contour certificate; several branches share the same
input observations.  Naive nested composition recomputes demands and may hide
an illegal adaptive oracle call.  A contract-DAG compiler separates two
passes:
\begin{enumerate}
\item a reverse pass determines every observation needed for the requested
output;
\item a forward pass executes each node once and records one bounded trace.
\end{enumerate}

\subsection{DAG model}

Let nodes be topologically ordered \(0,\ldots,N-1\).  Node \(v\) has parents
\(p_0(v),\ldots,p_{k_v-1}(v)<v\).  Its contract contains:
\begin{itemize}
\item an output query domain \(Q_v\);
\item for each output query \(q\), a parent-demand tuple
\(D_v(q)=(q_0,\ldots,q_{k_v-1})\);
\item a primitive-recursive evaluator \(E_v\);
\item a local validity predicate \(C_v\);
\item a local \(\Prf\,\PRAAx\) derivation of contract preservation.
\end{itemize}
If query domains are numerical precision levels, join is maximum.  For richer
finite queries, assume a primitive-recursive join \(q\sqcup q'\) with proofs
that it dominates both.  The displayed \(D_v\) is static: it depends on the
output query and fixed finite metadata, not on an as-yet uncomputed approximate
value.  Data-dependent algorithms must provide a certified static majorant, a
finite preliminary summary stage, or the joined demands of all unresolved
branches.  This condition is what makes the reverse pass genuinely prior to
the forward pass.

\subsection{Reverse demand algorithm}

Initialize every node demand to the least query \(\bot\), except the output
nodes, which receive the user's requested queries.  Scan nodes in reverse
topological order.  At node \(v\), compute \(D_v(q_v)\) and join each component
into the demand of the corresponding parent.  Because the graph is finite and
the scan length is \(N\), this is primitive recursive.

\begin{theorem}[Soundness of reverse demand]
\label{thm:demandsound}
Let \(q_v\) be the demands produced by the reverse pass.  If each source
presentation supplies a valid observation at its computed demand, then every
node \(v\) can be evaluated at demand \(q_v\), and every requested output
contract holds.
\end{theorem}

\begin{proof}
Induct on the topological order.  Source nodes satisfy their contracts by the
assumption.  For a non-source node \(v\), the reverse pass joined the demand
\(D_v(q_v)_j\) into each parent \(p_j(v)\).  Hence the observation produced for
that parent is at least as strong as the node's local requirement.  Restrict it
if necessary, apply the local contract derivation, and obtain a valid output at
\(q_v\).  Continue through the finite order.  The requested outputs have the
queries with which the reverse pass was initialized, so their contracts hold.
\end{proof}

\subsection{Fan-out and sharing}

Suppose a node feeds two consumers requiring precisions \(p\) and \(q\).  The
producer is evaluated once at \(p\sqcup q\).  Restriction maps supply the two
weaker observations.  This is more than an optimization: it ensures that both
consumers see coherent data from one finite observation rather than unrelated
witnesses.

For relational presentations, sharing does not force a single exact candidate.
The producer may return a cloud at the joined demand, and each consumer uses a
compatible restriction cloud.  A global viability pass selects a finite
compatible path only when needed.

\subsection{The global forward trace}

Encode in the trace state:
\begin{itemize}
\item the current node index;
\item the complete demand table;
\item the finite table of already computed observations;
\item one accumulated validity Boolean;
\item optional local certificate cursors.
\end{itemize}
At each step, the transition evaluates the next node from its parent entries
and updates the validity bit by conjunction with \(C_v\).  The invariant says:
\begin{enumerate}
\item every completed node has a well-formed observation at its assigned
demand;
\item every completed local contract is true;
\item untouched table entries retain the default value;
\item the demand table equals the result of the reverse pass.
\end{enumerate}
Each clause is a bounded Boolean fold.  The invariant is one equation and can
be fed to \cref{thm:tracereflection}.

\begin{theorem}[Certified DAG compilation]
A finite acyclic graph of proof-carrying contracted transducers with static or
certifiably majorized demands compiles to one trace system.  From the local
\(\Prf\,\PRAAx\) contract proofs, iPRA derives the global output contract at
every requested finite query.
\end{theorem}

\begin{proof}
The reverse-demand correctness theorem supplies the semantic invariant.
Booleanize its bounded table clauses.  The base proof verifies the initialized
tables.  The step proof is a finite case distinction on the current node,
followed by the corresponding local contract proof and table-update equations.
The final proof projects the accumulated output entries and their validity
bits.  Apply finite-trace reflection and universally quantify the requested
queries and source inputs.
\end{proof}

\subsection{Guarded branches}

A branch node may be compiled in three safe ways.
\begin{enumerate}
\item If the condition is a decidable finite Boolean, execute exactly one
branch.
\item If an approximate quantity is separated from the branch boundary by a
positive margin, refine until the finite guard certifies one side, then execute
that branch.
\item Without separation, execute or describe both finite branches and return
a relational union whose observable output is proved stable.
\end{enumerate}
An exact branch on the sign of an arbitrary real name is not permitted.

\subsection{Loops in a program graph}

A bounded loop is represented as one trace node whose local proof is itself
obtained from \cref{thm:tracereflection}.  Alternatively, the compiler can
flatten all loops into the global state.  The modular route creates smaller
local proof obligations; the flattened route permits stronger sharing.  Both
are conservative because the proof terms are composed by substitution and cut.

\subsection{Lean outline}

The supplied \code{CertifiedProgram.lean} sketches a syntax with assignments,
checks, sequential composition, bounded loops, and guarded or relational
branches.  Completion should proceed by:
\begin{enumerate}
\item replacing abbreviated semantic fields by formula builders and
\(\Prf\,\PRAAx\) derivations;
\item defining a topological list representation with a bounded well-formedness
checker;
\item implementing the reverse pass as a \(\PR\) fold;
\item proving the demand invariant by quantifier-free induction;
\item compiling the forward table update to a \code{PRTraceSystem};
\item invoking \code{traceReflection} once.
\end{enumerate}
No higher-order evaluator is needed: each certified program is a fixed Lean
value from which fixed \(\PR\) symbols are constructed.
